\documentclass[11pt]{article}
\usepackage[a4paper,margin=25mm]{geometry}
\usepackage{amsmath,amssymb,mathtools}
\usepackage{booktabs,longtable,array}
\usepackage[hidelinks]{hyperref}
\usepackage[T1]{fontenc}
\usepackage{lmodern}
\usepackage{microtype}
\setlength{\parskip}{0.6em}
\setlength{\parindent}{0pt}
\title{Mathematical Formulation of the Multi-Fidelity Air--Ground Coverage Simulator}
\author{ARSControl}
\date{}

\begin{document}
\maketitle

\begin{abstract}
This document is the typeset mathematical reference for the coupled multi-fidelity coverage simulator. It specifies the implemented scalar fields, sensor model, central autoregressive Gaussian process, controller densities, aerial and ground controllers, dynamics, event ordering, and evaluation metrics.
\end{abstract}

\section{Scope and notation}


The equations are marked conceptually as follows:


\begin{itemize}
\item \textbf{Model} describes the mathematical model or control objective.
\item \textbf{Implementation} gives the discrete equation evaluated by the current code.
\item \textbf{Compatibility path} describes retained legacy behavior that is not active when the central multi-fidelity estimator supplies the controllers.
\end{itemize}


The main symbols are:


\begin{longtable}{p{0.34\textwidth} p{0.56\textwidth}}
\toprule
Symbol & Meaning \\
\midrule
\endfirsthead
\toprule
Symbol & Meaning \\
\midrule
\endhead
$\mathbf q=(x,y)^\top$ & A point in the two-dimensional workspace $\Omega$ \\
$\mathbf p_i$ & Position of robot $i$ \\
$\theta_i$ & Heading of robot $i$ \\
$f_L,f_H$ & Latent LOW- and HIGH-fidelity scalar fields \\
$\delta$ & Discrepancy between the scaled LOW field and the HIGH field \\
$\rho$ & Fixed autoregressive scaling coefficient \\
$\mu_H,\sigma_H^2$ & Posterior mean and marginal variance of $f_H$ \\
$\phi$ & Nonnegative, normalized controller importance density \\
$w_j$ & Numerical quadrature weight associated with query point $\mathbf q_j$ \\
$V_{ij}$ & Indicator that query point $j$ belongs to robot $i$'s Voronoi cell \\
$\Delta t$ & Simulation or controller time step \\
$\mathrm{sat}_{[a,b]}(z)$ & Clipping of $z$ to the interval $[a,b]$ \\
\bottomrule
\end{longtable}


Unless stated otherwise, the current multi-fidelity path uses one central posterior for both teams:


\begin{itemize}
\item aerial robots use a normalized combination of posterior interest and uncertainty;
\item ground robots use posterior interest only;
\item ground positions affect aerial motion only indirectly, through the locations at which ground robots acquire HIGH-fidelity samples;
\item ground Voronoi cells are built from ground-robot positions only.
\end{itemize}


\section{Circular ground-obstacle map}


The coupled simulation maintains separate aerial and ground occupancy rasters. Let $M^A_{yx}\in\{0,1\}$ be the loaded aerial/HEDAC map. Robot positions and controller grids use raster-index coordinates independently of the HEDAC heat equation's physical cell-size parameter, so the production obstacle cell-centre coordinate is


\begin{equation}
\mathbf q_{yx}=
\begin{bmatrix}x&y\end{bmatrix}^{\!\top}.
\end{equation}


For episode seed $s$, circular-obstacle generation uses the independent deterministic seed $s_o=s+30000$. Given count $K_o$ and common radius $R_o$, candidate centres are sampled uniformly inside the map boundary inset by $R_o$ and accepted sequentially only if


\begin{equation}
\lVert\mathbf c_k-\mathbf c_l\rVert_2>2R_o
\quad\text{for every previously accepted }l.
\end{equation}


The exact generated occupancy and ground-map union are


\begin{equation}
O_{yx}
=
\mathbf 1\!\left[
\min_{k=1,\ldots,K_o}
\lVert\mathbf q_{yx}-\mathbf c_k\rVert_2
\le R_o
\right],
\qquad
M^G_{yx}=\max(M^A_{yx},O_{yx}).
\end{equation}


An infeasible request is rejected instead of silently reducing $K_o$. By
default, initial positions retain the historical uniform-free-cell sampling.
The composition experiment instead deploys both classes in the lower-left
corner. For class $r\in\{A,G\}$, map width $W$, height $H$, and $\gamma=0.2$,
the candidate set is


\begin{equation}
\mathcal C^r
=
\{(x,y):M^r_{yx}=0,\ 0\le x<\gamma W,\ 0\le y<\gamma H\}.
\end{equation}


Independent episode-derived seeds $s_A=s+40000$ and $s_G=s+50000$ generate
complete aerial and ground candidate permutations. Class $r$ receives the
first $N_r$ cells without replacement, followed by headings drawn from
$\mathcal U[0,2\pi)$. Consequently, the first class-specific initial states
are nested across compositions sharing an episode seed. An insufficient
corner candidate set is rejected rather than expanded silently. The HIGH truth
is zeroed where $M^G=1$, and nearest-cell samples of $M^G$ define the estimator
and evaluation query mask. Generated circles do not alter $M^A$: aerial
vehicles may overfly them.


\section{Simulator ground truth}


\subsection{HIGH-fidelity field}


The simulator constructs a smooth HIGH-fidelity field from a Gaussian mixture. For $K$ components,


\begin{equation}
g(\mathbf q) = \sum_{k=1}^{K}
\pi_k
\frac{
\exp\!\left[
-\frac{1}{2}
(\mathbf q-\boldsymbol\mu_k)^\top
\boldsymbol\Sigma_k^{-1}
(\mathbf q-\boldsymbol\mu_k)
\right]
}{
2\pi\sqrt{\det(\boldsymbol\Sigma_k)}
},
\sum_{k=1}^{K}\pi_k=1
\end{equation}


On the raster, this field is min-max scaled and obstacles are masked:


\begin{equation}
f_H(\mathbf q_j) \\
= \\
\mathbf 1_{\mathrm{free}}(\mathbf q_j) \\
\frac{g(\mathbf q_j)-g_{\min}} \\
{g_{\max}-g_{\min}+\varepsilon}
\end{equation}


This is the simulator's hidden truth. The estimator receives noisy samples, not the mixture parameters or the complete raster.


\subsection{Simulator LOW field and discrepancy}


The LOW truth is a normalized Gaussian smoothing of the masked HIGH truth. If $G_s$ is a Gaussian filter with standard deviation $s$ cells and $M$ is the free-space mask, the implemented normalized convolution is


\begin{equation}
f_L \\
= \\
\frac{G_s * (M f_H)}{G_s * M}
\end{equation}


where division is pointwise wherever the denominator is nonzero, and the result is zero outside the mask. The simulator then defines


\begin{equation}
\delta(\mathbf q)=f_H(\mathbf q)-\rho f_L(\mathbf q)
\end{equation}


so the autoregressive identity is exact for the simulated fields:


\begin{equation}
f_H(\mathbf q)=\rho f_L(\mathbf q)+\delta(\mathbf q)
\end{equation}


The simulator uses this identity to generate a coherent test problem. The estimator does not observe the simulator's $\delta$ field directly.


\section{Sensor model}


An observation from fidelity $s\in\{L,H\}$ at position $\mathbf q_n$ is


\begin{equation}
y_n=f_s(\mathbf q_n)+\epsilon_n, \qquad
\epsilon_n\sim\mathcal N(0,\tau_{s,n}^2)
\end{equation}


The configured sensor noise values are variances, so the noise standard deviation used to draw a sample is $\sqrt{\tau_{s,n}^2}$. Field values are read from the nearest raster cell after the continuous sample position is generated.


\subsection{Uniform-sector sampling}


For a robot at $\mathbf p$ with heading $\theta$, sensor range $R$, and total field-of-view angle $\Psi$, the sensor's only sampling law draws


\begin{equation}
\alpha\sim \\
\mathcal U\!\left(\theta-\frac{\Psi}{2}, \\
\theta+\frac{\Psi}{2}\right),
\qquad
U\sim\mathcal U(0,1),
\qquad
r=R\sqrt{U}
\end{equation}


and returns


\begin{equation}
\mathbf q \\
= \\
\mathbf p+ \\
r \\
\begin{bmatrix}
\cos\alpha\\
\sin\alpha \\
\end{bmatrix}
\end{equation}


The square root makes sample locations uniform with respect to area inside the sector because the polar area element is $r\,dr\,d\alpha$. Out-of-bounds draws are resampled up to the configured attempt limit. There is no sampling-mode parameter or alternative fixed-ray branch in the production sensor.


\section{Autoregressive multi-fidelity Gaussian process}


\subsection{Prior}


The central estimator assumes independent zero-mean Gaussian processes


\begin{equation}
f_L\sim\mathcal{GP}(0,k_L),
\qquad
\delta\sim\mathcal{GP}(0,k_\delta),
\qquad
f_L\perp\delta
\end{equation}


with


\begin{equation}
f_H(\mathbf q)=\rho f_L(\mathbf q)+\delta(\mathbf q)
\end{equation}


Both covariance functions are isotropic squared-exponential kernels:


\begin{equation}
k_a(\mathbf q,\mathbf q') \\
= \\
\sigma_a^2 \\
\exp\!\left( \\
-\frac{\lVert\mathbf q-\mathbf q'\rVert_2^2}{2\ell_a^2} \\
\right), \qquad
a\in\{L,\delta\}
\end{equation}


Here $\ell_a$ is a length scale and $\sigma_a^2$ is the configured kernel variance.


\subsection{What is known and what is learned}


\begin{longtable}{p{0.34\textwidth} p{0.56\textwidth}}
\toprule
Quantity & Current treatment \\
\midrule
\endfirsthead
\toprule
Quantity & Current treatment \\
\midrule
\endhead
Kernel family & Known: isotropic squared-exponential \\
Autoregressive structure & Known: $f_H=\rho f_L+\delta$ \\
$\rho$ & Fixed from configuration \\
Sensor noise variances & Fixed from configuration and stored per observation \\
LOW and HIGH sample positions and values & Observed \\
$\ell_L,\sigma_L^2,\ell_\delta,\sigma_\delta^2$ & Optionally fitted by marginal likelihood \\
Latent field values away from samples & Inferred as a posterior distribution \\
$\mu_H(\mathbf q),\sigma_H^2(\mathbf q)$ & Recomputed on the query grid after a successful update \\
Simulator truth, GMM parameters, and true discrepancy & Unknown to the estimator \\
\bottomrule
\end{longtable}


Fitting hyperparameters is empirical Bayes: it produces one bounded maximum-marginal-likelihood point estimate. It does not maintain a posterior distribution over the hyperparameters.


\subsection{Joint observation covariance}


Let


\begin{equation}
\mathbf y
=
\begin{bmatrix}
\mathbf y_L\\
\mathbf y_H
\end{bmatrix}
\end{equation}


with LOW inputs $X_L$, HIGH inputs $X_H$, and diagonal observation-noise matrices $D_L$ and $D_H$. In LOW-then-HIGH order, the training covariance is


\begin{equation}
K_y \\
= \\
\begin{bmatrix}
K_L(X_L,X_L)+D_L \\
\rho K_L(X_L,X_H) \\
\rho K_L(X_H,X_L) \\
\rho^2 K_L(X_H,X_H)+K_\delta(X_H,X_H)+D_H \\
\end{bmatrix}
\end{equation}


The implementation adds adaptive diagonal jitter:


\begin{equation}
\widetilde K_y=K_y+\eta I
\end{equation}


Starting from the configured $\eta$, failed Cholesky attempts multiply it by the configured jitter multiplier. No explicit matrix inverse is formed.


\subsection{Cholesky solve}


For


\begin{equation}
LL^\top=\widetilde K_y
\end{equation}


the implementation solves


\begin{equation}
L\mathbf z=\mathbf y, \\
\qquad \\
L^\top\boldsymbol\alpha=\mathbf z
\end{equation}


Thus


\begin{equation}
\boldsymbol\alpha=\widetilde K_y^{-1}\mathbf y
\end{equation}


conceptually, while the numerical computation uses triangular solves.


\subsection{HIGH-fidelity posterior}


For query points $X_*$, the covariance between training observations and the latent HIGH field is


\begin{equation}
K_{y*} = \begin{bmatrix}
\rho K_L(X_L,X_*) \\
\rho^2 K_L(X_H,X_*)+K_\delta(X_H,X_*) \\
\end{bmatrix}
\end{equation}


The prior HIGH covariance is


\begin{equation}
K_{**} \\
= \\
\rho^2K_L(X_*,X_*)+K_\delta(X_*,X_*)
\end{equation}


The posterior mean is


\begin{equation}
\boldsymbol\mu_H \\
= \\
K_{y*}^\top\boldsymbol\alpha
\end{equation}


If $V$ solves


\begin{equation}
LV=K_{y*}
\end{equation}


then the posterior marginal variance at query point $j$ is


\begin{equation}
\sigma_{H,j}^2 =
[K_{**}]_{jj}
-
\sum_r V_{rj}^2
\end{equation}


Only tiny negative values attributable to floating-point roundoff are clipped to zero. Materially negative variances cause the candidate update to fail.


\subsection{Hyperparameter fitting}


The four positive kernel parameters are optimized in log space:


\begin{equation}
\boldsymbol\xi \\
= \\
\log
\begin{bmatrix}
\ell_L \\
\sigma_L^2  \\
\ell_\delta \\
\sigma_\delta^2 \\
\end{bmatrix}^{\!\top}
\end{equation}


The bounded L-BFGS-B objective is the negative log marginal likelihood


\begin{equation}
\mathcal L(\boldsymbol\xi) \\
= \\
\frac{1}{2}\mathbf y^\top\boldsymbol\alpha \\
+ \\
\sum_{r=1}^{N}\log L_{rr} \\
+ \\
\frac{N}{2}\log(2\pi)
\end{equation}


Fitting occurs only when optimization is enabled, the retained sample count reaches its configured minimum, and the current posterior version satisfies the configured update interval. The value of $\rho$ and both sensor-noise models remain fixed.


\subsection{Query-grid quadrature}


For each axis, the estimator uses equally spaced query coordinates and trapezoidal weights. With spacing $h_x$,


\begin{equation}
w^x_j \\
= \\
\begin{cases}
h_x/2, & j\text{ is an endpoint},\\
h_x, & \text{otherwise}, \\
\end{cases}
\end{equation}


and similarly for $w^y_k$. The two-dimensional weight is


\begin{equation}
w_{jk}=w^x_jw^y_k
\end{equation}


These weights approximate integrals and are the $w_j$ used in density normalization and Lloyd centroid calculations.


\section{From the GP posterior to controller densities}


\subsection{Positive posterior interest}


The shared posterior interest field is the positive part of the HIGH posterior mean:


\begin{equation}
\mu_H^+(\mathbf q_j) \\
= \\
\max\!\left(\mu_H(\mathbf q_j),0\right)
\end{equation}


For a free-space mask $M_j\in\{0,1\}$ and quadrature weights $w_j$, define
$Z=\sum_k M_k\mu_H^+(\mathbf q_k)w_k$ and the startup-safe interest values


\begin{equation}
v_j
=
\begin{cases}
\mu_H^+(\mathbf q_j), & Z>0,\\
1, & Z=0.
\end{cases}
\end{equation}


The published density is


\begin{equation}
\phi_j \\
= \\
\frac{ \\
M_jv_j \\
}{ \\
\sum_k M_kv_kw_k \\
}
\end{equation}


Therefore
\begin{equation}
\phi_j\ge 0, \\
\qquad \\
\sum_j\phi_jw_j=1
\end{equation}


The clipping preserves all positive posterior contrast. It does not
exponentiate the field, shift negative values upward, or retain the removed
softplus transformation. If every posterior mean value is nonpositive, the
estimator publishes a uniform free-space density while retaining the raw mean
and variance in the snapshot. This permits a valid first posterior when robots
begin in a low-signal corner. Other prediction or normalization failures still
preserve the latest valid snapshot transactionally.


\subsection{Aerial interest-plus-uncertainty target}


Let


\begin{equation}
\bar\sigma_{H,j} \\
= \\
\begin{cases}
\sigma_{H,j}/\max_k\sigma_{H,k}, &
\max_k\sigma_{H,k}>0,\\
0, & \text{otherwise}. \\
\end{cases}
\end{equation}


The unnormalized aerial importance is


\begin{equation}
\widetilde\phi^{\,A}_j \\
= \\
\lambda_I\phi_j+\lambda_U\bar\sigma_{H,j}
\end{equation}


where $\lambda_I\ge 0$ is the interest weight and $\lambda_U\ge 0$ is the uncertainty weight. After masking and quadrature normalization,


\begin{equation}
\phi^A_j \\
= \\
\frac{M_j\widetilde\phi^{\,A}_j}{\sum_kM_k\widetilde\phi^{\,A}_kw_k}
\end{equation}


This density is bilinearly resampled onto the aerial HEDAC raster and normalized again there. All aerial robots see the same target; HEDAC coverage history and their different positions make their individual motions differ.


\subsection{Ground importance}


The ground controllers use


\begin{equation}
\phi^G_j=\phi_j
\end{equation}


Posterior variance is intentionally absent from the current ground objective. If the controller grid differs from the posterior grid, the density is bilinearly interpolated and quadrature-normalized again.


\section{HEDAC aerial coverage}


\subsection{Coverage footprint}


For an offset $\mathbf r$ from an aerial robot, the discrete coverage footprint is


\begin{equation}
\kappa(\mathbf r) \\
= \\
\exp\!\left( \\
-\frac{\lVert\mathbf r\rVert_2^2}{R_a} \\
\right)
\end{equation}


because the implementation uses the RBF shape parameter $1/R_a$. The square footprint is truncated using the configured minimum kernel value $\kappa_{\min}$. With workspace dimension $d=2$, its per-axis half-width is


\begin{equation}
h_\kappa \\
= \\
\left\lceil \\
\sqrt{ \\
\frac{-R_a\log\kappa_{\min}}{d} \\
} \\
\right\rceil
\end{equation}


At HEDAC step $n$, the cumulative raster coverage is updated by adding one footprint for every aerial robot:


\begin{equation}
C_j^{n+1} \\
= \\
C_j^n+\sum_i\kappa(\mathbf q_j-\mathbf p_i^n)
\end{equation}


with each footprint cropped at map boundaries.


\subsection{Coverage deficit and source}


HEDAC uses sum normalization on its raster:


\begin{equation}
\widehat C_j \\
= \\
\frac{M_j C_j}{\sum_kM_kC_k+\varepsilon}
\end{equation}


The current target $\phi^A$ is also sum-normalized on this raster. The positive coverage deficit and source are


\begin{equation}
d_j=\max(\phi^A_j-\widehat C_j,0), \\
\qquad \\
\widetilde S_j=d_j^2
\end{equation}


\begin{equation}
S_j \\
= \\
A_\Omega \\
\frac{M_j\widetilde S_j} \\
{\sum_kM_k\widetilde S_k+\varepsilon}
\end{equation}


where $A_\Omega$ is the free map area.


\subsection{Heat model}


The intended heat-equation form is


\begin{equation}
\frac{\partial T}{\partial t} \\
= \\
\alpha\nabla^2T \\
+ \\
sS \\
- \\
\beta T \\
- \\
\gamma L_c
\end{equation}


where $\alpha$ is diffusion, $s$ is source strength, $\beta$ is global cooling, and $L_c$ is local cooling.


The current non-optimized implementation evaluates the following five-point interior stencil exactly:


\begin{equation}
\mathcal D_\alpha[T]_{r,c} \\
= \\
\alpha \\
\left( \\
T_{r+1,c}+T_{r-1,c}+T_{r,c+1}+T_{r,c-1} \\
\right) \\
- \\
4T_{r,c}
\end{equation}


Then


\begin{equation}
T_{r,c}^{n+1} \\
= \\
T_{r,c}^{n} \\
+ \\
\Delta t_H \\
\left[ \\
\frac{\mathcal D_\alpha[T^n]_{r,c}}{\Delta x^2} \\
+ \\
sS_{r,c} \\
- \\
\frac{\beta}{A_\Omega}T_{r,c}^n \\
- \\
\frac{\gamma}{A_\Omega}(L_c)_{r,c} \\
\right]
\end{equation}


This discrete stencil is recorded exactly as implemented: the central $-4T_{r,c}$ term is not multiplied by $\alpha$. It therefore differs from the usual discretization $\alpha(T_{\mathrm{neighbors}}-4T_{r,c})/\Delta x^2$ when $\alpha\ne1$.


The active local-cooling raster is reset to zero every step and footprint accumulation into it is disabled, so the local-cooling contribution is currently zero. The non-optimized update modifies interior cells only, leaves outer boundary cells fixed, and does not explicitly skip obstacle cells during diffusion.


The explicit heat step is limited by


\begin{equation}
\Delta t_H \\
= \\
\min\!\left( \\
\Delta t, \\
\frac{c_{\mathrm{CFL}}\Delta x^2}{4\alpha} \\
\right)
\end{equation}


for $\alpha>0$, and $\Delta t_H=\Delta t$ for $\alpha=0$.


\subsection{Heat-gradient direction}


The raster gradients are computed by finite differences. They are first globally scaled by


\begin{equation}
g_{\mathrm{rms}} \\
= \\
\sqrt{ \\
\mathrm{mean}(T_x^2) \\
+ \\
\mathrm{mean}(T_y^2) \\
}
\end{equation}


when this value is nonzero. The scaled gradient is bilinearly interpolated at the robot position. Close to the outer boundary, the relevant component is replaced by a fixed inward value of magnitude $10$. Obstacle-wall avoidance is not active in this gradient function.


Finally, only the direction is retained:


\begin{equation}
\mathbf d_i \\
= \\
\begin{cases}
\nabla T(\mathbf p_i)/\lVert\nabla T(\mathbf p_i)\rVert_2, &
\lVert\nabla T(\mathbf p_i)\rVert_2>0,\\
\mathbf 0, & \text{otherwise}. \\
\end{cases}
\end{equation}


The target velocity and heading are


\begin{equation}
\mathbf v_i^\star=v_{\max}\mathbf d_i, \\
\qquad \\
\theta_i^\star=\mathrm{atan2}(d_{i,y},d_{i,x})
\end{equation}


\section{Robot dynamics}


\subsection{Aerial Dubins dynamics}


The canonical aerial configuration uses fixed-speed Dubins dynamics:


\begin{equation}
\dot x_i=v_A\cos\theta_i, \\
\qquad \\
\dot y_i=v_A\sin\theta_i, \\
\qquad \\
\dot\theta_i=u_i
\end{equation}


The coordinated-turn limit derived from maximum bank angle $\varphi_{\max}$ is


\begin{equation}
u_{\max} \\
= \\
\frac{g\tan\varphi_{\max}}{v_A}
\end{equation}


The HEDAC heading tracker uses


\begin{equation}
e_{\theta,i} \\
= \\
\mathrm{atan2} \\
\left( \\
\sin(\theta_i^\star-\theta_i), \\
\cos(\theta_i^\star-\theta_i) \\
\right)
\end{equation}


\begin{equation}
u_i \\
= \\
\mathrm{sat}_{[-u_{\max},u_{\max}]} \\
\left(2e_{\theta,i}\right)
\end{equation}


The implementation updates heading first and then position:


\begin{equation}
\theta_i^{n+1} \\
= \\
\mathrm{wrap} \\
\left(\theta_i^n+u_i^n\Delta t_A\right)
\end{equation}


\begin{equation}
\mathbf p_i^{n+1} \\
= \\
\mathbf p_i^n \\
+ \\
v_A\Delta t_A \\
\begin{bmatrix} \\
\cos\theta_i^{n+1}\\ \\
\sin\theta_i^{n+1} \\
\end{bmatrix}
\end{equation}


Positions are clipped to the map after the controller step.


\subsection{Ground unicycle dynamics}


Both current ground controllers require or predict the unicycle model


\begin{equation}
\dot x_i=v_i\cos\theta_i, \\
\qquad \\
\dot y_i=v_i\sin\theta_i, \\
\qquad \\
\dot\theta_i=\omega_i
\end{equation}


The real ground agent clips $(v_i,\omega_i)$ to its configured limits, updates heading first, and then position:


\begin{equation}
\theta_i^{n+1} \\
= \\
\mathrm{wrap} \\
\left(\theta_i^n+\omega_i^n\Delta t_G\right)
\end{equation}


When the ground map contains obstacles, the coupled controller predicts this heading-first segment before committing translation. The segment is sampled with spacing no greater than $0.25$ map units. With $\mathcal S_i^n$ denoting those samples, the executed velocity is


\begin{equation}
v_{i,\mathrm{exec}}^n
=
\begin{cases}
v_i^n,
& \mathcal S_i^n\text{ lies inside the raster and }M^G(\mathbf q)=0
\text{ for all }\mathbf q\in\mathcal S_i^n,\\
0,&\text{otherwise}.
\end{cases}
\end{equation}


The clipped angular command is applied in both cases, allowing a blocked robot to turn away. When $M^G$ contains no occupied cell the guard is bypassed exactly, preserving the previous obstacle-free state update. This guard is used by Lloyd control and by the compatibility MPC path when its agent is a unicycle.


\begin{equation}
\mathbf p_i^{n+1} \\
= \\
\mathbf p_i^n \\
+ \\
v_i^n\Delta t_G \\
\begin{bmatrix} \\
\cos\theta_i^{n+1}\\ \\
\sin\theta_i^{n+1} \\
\end{bmatrix}
\end{equation}


The MPC prediction model uses forward Euler with the old heading instead:


\begin{equation}
\begin{bmatrix} \\
x_{h+1}\\y_{h+1}\\\theta_{h+1} \\
\end{bmatrix} \\
= \\
\begin{bmatrix} \\
x_h+v_h\cos\theta_h\,\Delta t\\ \\
y_h+v_h\sin\theta_h\,\Delta t\\ \\
\theta_h+\omega_h\Delta t \\
\end{bmatrix}
\end{equation}


This one-step discretization difference is part of the present implementation.


\section{Ground-only Voronoi partition}


For ground positions $\{\mathbf p_1,\ldots,\mathbf p_R\}$, query point $\mathbf q_j$ is assigned by


\begin{equation}
i^\star(j) \\
= \\
\operatorname*{arg\,min}_{i\in\{1,\ldots,R\}} \\
\lVert\mathbf q_j-\mathbf p_i\rVert_2^2
\end{equation}


\begin{equation}
V_{ij} \\
= \\
\begin{cases}
1, & i=i^\star(j)
\text{ and }
\lVert\mathbf q_j-\mathbf p_i\rVert_2\le R_V,\\ \\
0, & \text{otherwise}. \\
\end{cases}
\end{equation}


Only ground positions are passed to this partition. Aerial positions do not compete for ground Voronoi cells. The Lloyd controller chooses a range larger than the map diagonal, so its cells cover the entire query grid. The MPC controller currently uses $R_V=50$ in map-coordinate units.


\section{Lloyd ground controller}


\subsection{Continuous coverage objective}


The classical locational coverage objective is


\begin{equation}
\mathcal H(\mathbf p_1,\ldots,\mathbf p_R) \\
= \\
\sum_{i=1}^{R} \\
\int_{\mathcal V_i} \\
\lVert\mathbf q-\mathbf p_i\rVert_2^2 \\
\phi^G(\mathbf q)\,d\mathbf q
\end{equation}


For fixed Voronoi cells, its minimizer places each generator at the density-weighted centroid of its cell.


\subsection{Discrete mass and centroid}


The implemented cell mass is


\begin{equation}
m_i \\
= \\
\sum_j \\
V_{ij}\phi^G_jw_j
\end{equation}


The centroid is


\begin{equation}
\mathbf c_i \\
= \\
\frac{ \\
\sum_j \\
\mathbf q_jV_{ij}\phi^G_jw_j \\
}{ \\
m_i \\
}
\end{equation}


Thus $w_j$ is the area represented by grid point $j$ under trapezoidal quadrature. On a uniform interior grid it is approximately $\Delta x\Delta y$; boundary weights are halved per boundary axis. These weights are what make the vectorized sum approximate the slower nested area-integral calculation.


If $m_i$ is numerically zero, the current robot position is used as the centroid.


\subsection{Centroid tracking}


Define the raw centroid displacement


\begin{equation}
\mathbf e_i^c=\mathbf c_i-\mathbf p_i.
\end{equation}


For obstacle maps, the Lloyd path first augments this displacement. Let $\mathcal O_i$ contain occupied cell centres within


\begin{equation}
R_I=\max(2,2.5R_o)
\end{equation}


of robot $i$, and let $d_{il}=\lVert\mathbf p_i-\mathbf o_l\rVert_2$. The bounded repulsion is


\begin{equation}
\mathbf r_i
=
\operatorname{clip}_{\lVert\cdot\rVert\le1}
\left[
\sum_{l\in\mathcal O_i}
\frac{R_I-d_{il}}{R_I}
\frac{\mathbf p_i-\mathbf o_l}{d_{il}}
\right].
\end{equation}


The displacement used by the heading and speed law is then


\begin{equation}
\widetilde{\mathbf e}_i
=
\mathbf c_i-\mathbf p_i
+w_o\max(\lVert\mathbf c_i-\mathbf p_i\rVert_2,1)\mathbf r_i,
\end{equation}


where $w_o$ is the configured ground wall-avoidance weight. The tracking displacement is $\mathbf e_i=\mathbf e_i^c$ on an obstacle-free map and $\mathbf e_i=\widetilde{\mathbf e}_i$ otherwise. Define


\begin{equation}
r_i=\lVert\mathbf e_i\rVert_2,
\qquad
\theta_i^\star=\mathrm{atan2}(e_{i,y},e_{i,x}).
\end{equation}


\begin{equation}
e_{\theta,i} \\
= \\
\mathrm{wrap}(\theta_i^\star-\theta_i)
\end{equation}


Outside the centroid tolerance, the controls are


\begin{equation}
v_i \\
= \\
\mathrm{sat}_{[-v_{\max},v_{\max}]} \\
\left(k_p r_i\cos e_{\theta,i}\right)
\end{equation}


\begin{equation}
\omega_i \\
= \\
\mathrm{sat}_{[-\omega_{\max},\omega_{\max}]} \\
\left(k_\theta e_{\theta,i}\right)
\end{equation}


Inside the tolerance, both controls are zero. The factor $\cos e_{\theta,i}$ reduces forward motion when the centroid is lateral and allows bounded reverse motion when it lies behind the robot.


\section{MPC ground controller}


\subsection{Per-robot importance weights}


Robot $i$ receives the masked weights


\begin{equation}
W_{ij}=V_{ij}\phi^G_j
\end{equation}


These weights do not include quadrature factors in the current MPC objective.


\subsection{Implemented limited-FOV stage cost}


For predicted state $\mathbf x_h=(x_h,y_h,\theta_h)^\top$, define


\begin{equation}
r_{hj}^2 \\
= \\
\left\lVert \\
\mathbf q_j- \\
\begin{bmatrix}x_h\\y_h\end{bmatrix} \\
\right\rVert_2^2
\end{equation}


\begin{equation}
a_{hj} \\
= \\
\mathrm{atan2} \\
(q_{j,y}-y_h,q_{j,x}-x_h)-\theta_h
\end{equation}


\begin{equation}
F_{hj}=2-\frac{r_{hj}^2}{R_F^2}, \\
\qquad \\
M_{hj} \\
= \\
\exp\!\left[ \\
-\frac{(a_{hj}-\Psi/2)^2}{2(\Psi/2)^2} \\
\right]
\end{equation}


The exact current stage cost is


\begin{equation}
\ell_i(\mathbf x_h) \\
= \\
- \\
\sum_jF_{hj}M_{hj}W_{ij} \\
- \\
3\sum_j \\
W_{ij} \\
\exp\!\left( \\
-\frac{r_{hj}^2}{3R_F^2} \\
\right)
\end{equation}


This formula is documented as implemented. In particular, the angular Gaussian is centered at $+\Psi/2$, not at zero, and the relative angle is not explicitly wrapped. The function parameter intended for mask steepness is currently unused.


\subsection{Finite-horizon problem}


With horizon $H$, the solver computes


\begin{equation}
\min_{\{v_h,\omega_h\}_{h=0}^{H-1}} \\
\quad \\
10\sum_{h=0}^{H-1}\ell_i(\mathbf x_h
\end{equation}


subject to the unicycle prediction dynamics and control bounds


\begin{equation}
-a_{\max}\le v_h\le a_{\max}, \\
\qquad \\
-a_{\max}\le\omega_h\le a_{\max}
\end{equation}


The same configured quantity named maximum acceleration is currently used as the bound for both MPC control components. Only the terminal predicted position is constrained:


\begin{equation}
0\le x_H\le W_{\mathrm{map}}, \\
\qquad \\
0\le y_H\le H_{\mathrm{map}}
\end{equation}


There is currently no active collision penalty, control-effort penalty, or orientation penalty in this optimization. IPOPT is warm-started from the previous solution, and only the first optimized pair $(v_0,\omega_0)$ is sent to the real unicycle agent.


\subsection{Other retained cost functions}


The codebase also defines, but the current coupled MPC does not use,


\begin{equation}
J_{\mathrm{coverage}}(\mathbf p_i) \\
= \\
\sum_j \\
\lVert\mathbf q_j-\mathbf p_i\rVert_2^2W_{ij}
\end{equation}


and the collision penalty


\begin{equation}
J_{\mathrm{collision}} \\
= \\
\alpha_c \\
\exp\!\left[ \\
-\beta_c \\
\left( \\
\lVert\mathbf p_i-\mathbf p_{\mathrm{obs}}\rVert_2^2-D_s^2 \\
\right) \\
\right]
\end{equation}


They are included here to distinguish available mathematical utilities from the active MPC objective.


\section{Asynchronous estimation and causal order}


For a periodic event with period $P$, start time $t_0$, and fire count $k$, the next scheduled time is computed without accumulated drift:


\begin{equation}
t_{\mathrm{next}}=t_0+kP
\end{equation}


An event is due when


\begin{equation}
t+\varepsilon_t\ge t_{\mathrm{next}}
\end{equation}


At simulation step $n$, $t_n=n\Delta t$. The current coupled order is:


\begin{enumerate}
\item collect LOW samples from aerial robots if due;
\item read the latest already-published posterior and move aerial robots with HEDAC;
\item collect HIGH samples from ground robots if due;
\item update and publish the central GP if due;
\item read that posterior density and move ground robots.
\end{enumerate}


Consequently, a HIGH sample collected at step $n$ can influence ground control at step $n$ after a successful update, but influences aerial control no earlier than step $n+1$.


When video recording is enabled, it is an observer of this completed state:
after each outer step, it may render the already-published posterior and the
current aerial and ground positions. It captures only steps satisfying
$n \bmod r=0$, where $r$ is the configured frame interval, and also captures
the final completed step if it was not already selected. Rendering and encoding
do not submit observations, update the estimator, or alter either controller.
For one video, the uncertainty panel uses the fixed color interval
$[0,\sigma_{\max}^{(0)}]$, where $\sigma_{\max}^{(0)}$ is the largest HIGH
posterior standard deviation in its first captured frame; later frames reuse
that interval rather than independently rescaling their colors.


Estimator publication is transactional. A candidate update must complete retention, optional hyperparameter fitting, Cholesky factorization, prediction, positive-part conversion, and density validation before it replaces the cached posterior.


\section{Evaluation quantities}


\subsection{Current HEDAC coverage error}


The quantity reported by HEDAC as an ergodic metric is a discrete spatial $L^2$ difference. With sum-normalized cumulative coverage


\begin{equation}
\widehat C_j \\
= \\
\frac{C_j}{\sum_kC_k+\varepsilon}
\end{equation}


the reported value is


\begin{equation}
E \\
= \\
\sqrt{ \\
\sum_j \\
\left[ \\
M_j(\widehat C_j-\phi_j^{\mathrm{reference}}) \\
\right]^2 \\
}
\end{equation}


This is not a Fourier-coefficient ergodic metric. In the present HEDAC step, $\phi^{\mathrm{reference}}$ is the original normalized goal field stored when HEDAC was constructed, not necessarily the latest external multi-fidelity aerial target. It should therefore be interpreted carefully in coupled runs.


\subsection{Posterior reconstruction diagnostics}


Useful field diagnostics include root mean squared error


\begin{equation}
\mathrm{RMSE} \\
= \\
\sqrt{ \\
\frac{1}{N} \\
\sum_{j=1}^{N} \\
\left( \\
\mu_H(\mathbf q_j)-f_H(\mathbf q_j) \\
\right)^2 \\
}
\end{equation}


and mean negative log predictive density, when evaluating against simulator truth:


\begin{equation}
\mathrm{MNLPD} \\
= \\
\frac{1}{N} \\
\sum_{j=1}^{N} \\
\left[ \\
\frac{1}{2}\log(2\pi\sigma_{H,j}^2) \\
+ \\
\frac{ \\
\left(f_H(\mathbf q_j)-\mu_H(\mathbf q_j)\right)^2 \\
}{ \\
2\sigma_{H,j}^2 \\
} \\
\right]
\end{equation}


These diagnostics use hidden truth only for evaluation; they are not controller inputs.


\subsection{Closed-loop team-composition reconstruction evaluation}


The composition study fixes $N=N_a+N_g$ and varies the aerial and ground
counts while running the same coupled controllers and autoregressive MFGP. For
the publication study, $N=10$ and $N_a\in\{10,2,0\}$. The named scenario
parameters are


\begin{equation}
(K_s,T_s,n_{o,s},r_{o,s})
=
\begin{cases}
(200,20,5,1), & s=\texttt{easy\_long},\\
(100,10,15,2), & s=\texttt{hard\_short},
\end{cases}
\end{equation}


where $K_s$ is the number of integration steps, $T_s=K_s\Delta t$ seconds,
$n_{o,s}$ is the ground-obstacle count, $r_{o,s}$ is the obstacle radius, and
$\Delta t=0.1$ s. Geometry and duration both change, so this is a comparison
of bundled operating regimes rather than an independent-factor experiment. For
composition $c$, episode $e$, and posterior publication $r$, let
$\mathcal J_{c,e}$ denote the free query cells and define
$\mu^+_{H,c,e,r,j}=\max(\mu_{H,c,e,r,j},0)$. The implemented normalized
nonnegative-field reconstruction error is


\begin{equation}
\operatorname{NRMSE}_{c,e,r}
=
\frac{
\sqrt{
\frac{1}{|\mathcal J_{c,e}|}
\sum_{j\in\mathcal J_{c,e}}
(\mu^+_{H,c,e,r,j}-f_{H,c,e,j})^2
}
}{
\max_{j\in\mathcal J_{c,e}}f_{H,c,e,j}
-
\min_{j\in\mathcal J_{c,e}}f_{H,c,e,j}
}.
\end{equation}


With integration weights $w_j$, normalized hidden truth density
$\phi_{c,e,j}$, and the positive-floor-renormalized posterior density
$\widehat\phi_{c,e,r,j}$, the implemented KL divergence is


\begin{equation}
D_{\mathrm{KL},c,e,r}
=
\sum_{j:\phi_{c,e,j}>0}
w_j\phi_{c,e,j}
\log\!\left(
\frac{\phi_{c,e,j}}{\widehat\phi_{c,e,r,j}}
\right).
\end{equation}


Define the variance protected against numerical collapse by
$\widetilde\sigma^2_{H,c,e,r,j}
=\max(\sigma^2_{H,c,e,r,j},10^{-12})$. The implemented
integration-weighted marginal Gaussian negative log predictive density is


\begin{equation}
\operatorname{NLPD}_{c,e,r}
=
\frac{1}{\sum_{j\in\mathcal J_{c,e}}w_j}
\sum_{j\in\mathcal J_{c,e}}w_j
\left[
\frac{1}{2}\log(2\pi\widetilde\sigma^2_{H,c,e,r,j})
+
\frac{(f_{H,c,e,j}-\mu_{H,c,e,r,j})^2}
{2\widetilde\sigma^2_{H,c,e,r,j}}
\right].
\end{equation}


This score uses the raw latent HIGH posterior mean and variance. It applies no
nonnegative clipping or density normalization, adds no observation noise, and
uses the marginal predictive density at each query rather than the joint
cross-query covariance. Lower is better; a continuous predictive density can
produce a negative NLPD.


No density normalization is applied to either field score. The secondary
empirical calibration diagnostic remains centred on the raw Gaussian-process
mean rather than $\mu_H^+$:


\begin{equation}
C_{95,c,e,r}
=
\frac{1}{|\mathcal J_{c,e}|}
\sum_{j\in\mathcal J_{c,e}}
\mathbf 1\!\left[
|f_{H,c,e,j}-\mu_{H,c,e,r,j}|
\le
1.96\sqrt{\sigma^2_{H,c,e,r,j}}
\right].
\end{equation}


For a metric $m_r$ published at $t_0<\cdots<t_R\le T_s$, the evaluator holds
the latest value to the fixed mission horizon and computes


\begin{equation}
\overline m
=
\frac{1}{T_s-t_0}
\left[
\sum_{r=0}^{R-1}
\frac{m_r+m_{r+1}}{2}(t_{r+1}-t_r)
+m_R(T_s-t_R)
\right].
\end{equation}


One active retention budget $B$ is split in proportion to composition. For
$N_a,N_g>0$, the exact caps are


\begin{equation}
B_L
=
\left\lfloor B\frac{N_a}{N}\right\rfloor,
\qquad
B_H
=
B-B_L.
\end{equation}


At a homogeneous endpoint the active fidelity receives $B$; the inactive
buffer receives an unused capacity of one because the retention class requires
a positive cap. Both fidelities have the same sensing period and samples per
robot per event, so each condition has $sN$ observation opportunities per
sensing event. These budgets and all evaluation metrics are diagnostics and do
not alter the controller laws.


The composition experiment enables online empirical-Bayes kernel fitting.
Every composition and episode starts from the same configured parameter vector
$\boldsymbol\theta_0=(\ell_L,\sigma_L^2,\ell_\delta,\sigma_\delta^2)$ and uses
the same bounds, minimum sample count $M_{\min}$, fit interval
$K_{\mathrm{fit}}$, restart count, and iteration limit. For the candidate that
would publish posterior version $r\ge1$, the exact scheduling rule is


\begin{equation}
\boldsymbol\theta_{c,e,r}
=
\begin{cases}
\displaystyle\arg\min_{\boldsymbol\theta\in\mathcal B}
\mathcal L_{c,e,r}(\boldsymbol\theta),
& n_{L,c,e,r}+n_{H,c,e,r}\ge M_{\min}
\ \text{and}\ (r-1)\bmod K_{\mathrm{fit}}=0,\\
\boldsymbol\theta_{c,e,r-1}, & \text{otherwise}.
\end{cases}
\end{equation}


The policy is controlled across conditions, but the realized fitted values are
data-dependent and may differ. At A10/G0 the discrepancy parameters are not
identified by LOW observations. At A0/G10, HIGH observations identify the
combined HIGH covariance but do not separately identify its LOW and
discrepancy components. Endpoint fitted values are therefore optimizer
diagnostics rather than independent physical parameter estimates. Each saved
posterior record includes the fit flag, optimization duration, and all four
realized kernel parameters.


For the publication protocol, $B=400$ and the LOW/HIGH caps are $(400,1)$,
$(80,320)$, and $(1,400)$ for A10/G0, A2/G8, and A0/G10. With ten samples per
robot every $0.5$ s, each composition submits 100 observations per event:
4000 over the easy/long mission and 2000 over the hard/short mission.


\subsection{Footprint-normalized ground coverage effectiveness}


Let $\phi_j$ denote the hidden normalized truth density at query point $j$,
with integration weight $w_j$. The raw probability mass visible through the
union of the instantaneous ground sensing sectors is


\begin{equation}
M_{\mathrm{vis}}(t)
=
\sum_{j:\,\mathbf q_j\in\cup_i F_i(t)} \phi_j w_j.
\end{equation}


For $N_g$ ground robots, sensing radius $R$, and field-of-view angle $\theta$
in radians, one ideal sector has area $\theta R^2/2$. The fixed episode area
budget is


\begin{equation}
A_B
=
\min\left(
A_{\mathrm{free}},
N_g\frac{\theta R^2}{2}
\right),
\qquad
A_{\mathrm{free}}=\sum_{j\in\mathcal Q_{\mathrm{free}}}w_j.
\end{equation}


The discrete densest-area oracle is


\begin{equation}
M_\star(A_B)
=
\max_{0\le z_j\le 1}
\sum_{j\in\mathcal Q_{\mathrm{free}}}z_j\phi_jw_j,
\qquad
\sum_{j\in\mathcal Q_{\mathrm{free}}}z_jw_j\le A_B.
\end{equation}


The implementation solves this fractional selection exactly by sorting free
cells in decreasing order of $\phi_j$, accepting whole cells until the budget
is reached, and accepting the necessary fraction of the final cell. The
reported score is


\begin{equation}
E_{\mathrm{cov}}(t)
=
\frac{M_{\mathrm{vis}}(t)}{M_\star(A_B)}.
\end{equation}


The denominator is fixed within an episode. It relaxes sector geometry and is
therefore an optimistic upper bound for the same nominal sensing area. Actual
footprint overlap, boundary clipping, obstacle clipping, and poor placement
reduce the numerator. The evaluator saves both $M_{\mathrm{vis}}$ and
$E_{\mathrm{cov}}$; only hidden truth is used, and neither quantity feeds back
to a controller.

For paired proposed-method/Egerstedt evaluation, method $m$ uses its saved
ground states, field-of-view angle, and sensing range. At saved state time
$t_k$, the exact implemented numerator for episode $e$ is

\begin{equation}
M_{\mathrm{vis},k}^{(m,e)}
=
\sum_j w_j\phi_j^{(e)}
\mathbf 1\!\left[
\mathbf q_j\in\bigcup_iF_{i,k}^{(m,e)}
\right].
\end{equation}

The proposed method uses its saved headings to orient sectors. The Egerstedt
baseline declares a $360$-degree disk, so its absent heading is represented by
an immaterial zero. Each method receives its own equal-area oracle because the
saved footprint geometries can differ. The evaluator saves the final score and
the trapezoidal time average

\begin{equation}
\overline E_{\mathrm{cov}}^{(m,e)}
=
\frac{1}{t_K-t_0}
\sum_{k=0}^{K-1}
\frac{
E_{\mathrm{cov},k}^{(m,e)}
+
E_{\mathrm{cov},k+1}^{(m,e)}
}{2}
(t_{k+1}-t_k).
\end{equation}

Runtime comparison uses the saved total wall-clock duration of each complete
method step. The evaluated archive retains the raw step totals and their
per-episode mean, median, empirical 95th percentile, and maximum. These
diagnostics do not feed back into either controller.


\subsection{Separate Egerstedt comparison baseline}


The comparison baseline uses the saved obstacle-free scenario but calls only
the control functions in \texttt{evaluation/egerstedt.py}. It assumes the
hidden truth $\phi_j$ is known. For aerial robot $a$, let $V^A_a$ be its
range-limited Euclidean Voronoi set and let $c^A_a$ be the unweighted centroid
of its query-grid points. The implemented discrete aerial update is


\begin{equation}
p^A_{a,k+1}
=
p^A_{a,k}
+
\Delta t\,0.5\left(c^A_{a,k}-p^A_{a,k}\right).
\end{equation}


Each ground robot computes the truth-density-weighted centroid $c^G_i$ of its
own range-limited Voronoi set. Ground robots are assigned to their nearest
aerial robot. If $n_a$ is the number assigned to aerial cell $a$, the code
computes


\begin{equation}
\sigma_a
=
\frac{n_a}{N_g}
-
\frac{\sum_{j\in V^A_a}\phi_j}{\sum_j\phi_j},
\qquad
\widehat\sigma_a
=
\begin{cases}
\sigma_a,&\sigma_a>1/N_g,\\
0,&\text{otherwise}.
\end{cases}
\end{equation}


Let $a^- = \arg\min_a\sigma_a$ and let $a(i)$ be the aerial cell containing
ground robot $i$. With both baseline gains equal to one, the implemented
ground update is


\begin{equation}
p^G_{i,k+1}
=
p^G_{i,k}
+
\Delta t
\left[
\left(1-\widehat\sigma_{a(i)}\right)
\left(c^G_{i,k}-p^G_{i,k}\right)
+
\widehat\sigma_{a(i)}
\left(c^A_{a^-,k}-p^G_{i,k}\right)
\right].
\end{equation}


The runner uses the paired mission timestep $\Delta t$ rather than the
demonstration script's default timestep, clips ground positions to the
rectangular domain, and rejects maps containing obstacles. Both classes are
first-order point robots with omnidirectional range-limited sensing; no heading
state is fabricated. The stored baseline diagnostic is


\begin{equation}
H_k
=
\Delta A
\sum_j
\phi_j
\left\|\mathbf q_j-p^G_{\operatorname{Vor}(j),k}\right\|^2.
\end{equation}


This locational cost is not substituted for the common physical coverage
metric. Because the baseline consumes oracle truth directly, it publishes no
reconstructed density or posterior uncertainty.


\section{Legacy compatibility path}


When no central multi-fidelity posterior is supplied, the retained ground path can fuse separate aerial and ground GP estimates. Let $\bar\sigma_A,\bar\sigma_G$ be each standard-deviation field divided by its own maximum. The implemented weights are


\begin{equation}
d_j \\
= \\
\frac{1}{\bar\sigma_{G,j}+\varepsilon} \\
+ \\
\frac{1}{\bar\sigma_{A,j}+\varepsilon}
\end{equation}


\begin{equation}
\lambda_{A,j} \\
= \\
\frac{ \\
1/(\bar\sigma_{A,j}+\varepsilon) \\
}{d_j}, \\
\qquad \\
\lambda_{G,j} \\
= \\
\frac{ \\
1/(\bar\sigma_{G,j}+\varepsilon) \\
}{d_j}
\end{equation}


and


\begin{equation}
\mu_{\mathrm{fused},j} \\
= \\
\lambda_{A,j}\mu_{A,j} \\
+ \\
\lambda_{G,j}\mu_{G,j}
\end{equation}


This is a compatibility fallback, not the central autoregressive GP. In multi-fidelity mode the controllers consume the shared posterior products defined in Section 5.


\section{Equation-to-code map}


\begin{longtable}{p{0.34\textwidth} p{0.56\textwidth}}
\toprule
Mathematical component & Primary implementation \\
\midrule
\endfirsthead
\toprule
Mathematical component & Primary implementation \\
\midrule
\endhead
HIGH truth and Gaussian mixture & \href{../src/core/gmm.py}{\texttt{\detokenize{src/core/gmm.py}}}, \href{../src/coupled_simulation.py}{\texttt{\detokenize{src/coupled_simulation.py}}} \\
Circular obstacle generation, masking, corner initialization, and motion guard & \href{../src/core/obstacles.py}{\texttt{\detokenize{src/core/obstacles.py}}}, \href{../src/coupled_simulation.py}{\texttt{\detokenize{src/coupled_simulation.py}}} \\
LOW smoothing and sensor observations & \href{../src/models/sensors.py}{\texttt{\detokenize{src/models/sensors.py}}} \\
Autoregressive GP and marginal likelihood & \href{../src/core/multifidelity_gp.py}{\texttt{\detokenize{src/core/multifidelity_gp.py}}} \\
Asynchronous estimator and posterior publication & \href{../src/core/multifidelity_estimator.py}{\texttt{\detokenize{src/core/multifidelity_estimator.py}}} \\
Density clipping, normalization, and aerial target & \href{../src/core/density.py}{\texttt{\detokenize{src/core/density.py}}} \\
Event schedule and controller-density caching & \href{../src/simulation.py}{\texttt{\detokenize{src/simulation.py}}} \\
Video observer and encoding & \href{../src/utils/multifidelity_video.py}{\texttt{\detokenize{src/utils/multifidelity_video.py}}}, \href{../examples/run_multifidelity.py}{\texttt{\detokenize{examples/run_multifidelity.py}}} \\
HEDAC source, heat update, and gradient law & \href{../src/core/hedac.py}{\texttt{\detokenize{src/core/hedac.py}}}, \href{../src/utils/math_utils.py}{\texttt{\detokenize{src/utils/math_utils.py}}} \\
Dubins and unicycle state updates & \href{../src/models/agents.py}{\texttt{\detokenize{src/models/agents.py}}} \\
MPC prediction dynamics & \href{../src/models/models.py}{\texttt{\detokenize{src/models/models.py}}} \\
Voronoi partition and Lloyd centroids & \href{../src/utils/voronoi.py}{\texttt{\detokenize{src/utils/voronoi.py}}}, \href{../src/coupled_simulation.py}{\texttt{\detokenize{src/coupled_simulation.py}}} \\
MPC objective & \href{../src/core/costFunctions.py}{\texttt{\detokenize{src/core/costFunctions.py}}}, \href{../src/coupled_simulation.py}{\texttt{\detokenize{src/coupled_simulation.py}}} \\
Coverage error & \href{../src/core/base.py}{\texttt{\detokenize{src/core/base.py}}} \\
Composition-sweep reconstruction evaluation & \href{../evaluation/run_multifidelity_composition.py}{\texttt{composition runner}}, \href{../evaluation/multifidelity_metrics.py}{\texttt{metric equations}}, \href{../evaluation/evaluate_multifidelity_composition.py}{\texttt{offline evaluator}}, \href{../evaluation/plot_multifidelity_composition.py}{\texttt{standalone plotter}}, \href{../evaluation/run_multifidelity_scenario_pipeline.py}{\texttt{scenario pipeline}} \\
Footprint-normalized evaluation coverage & \href{../evaluation/multifidelity_metrics.py}{\texttt{\detokenize{evaluation/multifidelity_metrics.py}}}, \href{../evaluation/evaluate_multifidelity_ablation.py}{\texttt{\detokenize{evaluation/evaluate_multifidelity_ablation.py}}}, \href{../evaluation/evaluate_baseline_comparison.py}{\texttt{\detokenize{evaluation/evaluate_baseline_comparison.py}}} \\
Egerstedt comparison baseline & \href{../evaluation/egerstedt.py}{\texttt{\detokenize{evaluation/egerstedt.py}}}, \href{../evaluation/run_egerstedt_comparison.py}{\texttt{\detokenize{evaluation/run_egerstedt_comparison.py}}} \\
\bottomrule
\end{longtable}


The associated configuration parameters and their tuning effects are documented in \href{multifidelity_config_reference.md}{Multi-fidelity configuration reference}.


\end{document}
